\documentclass[12pt,a4paper]{article}

\usepackage{iftex}
\ifTUTeX % tectonic / xelatex / lualatex
  \usepackage{fontspec}
\else % pdflatex
  \usepackage[utf8]{inputenc}
  \usepackage[T1]{fontenc}
  \usepackage{lmodern}
\fi
\usepackage[polish]{babel}
\usepackage[a4paper,margin=2.5cm]{geometry}
\usepackage{graphicx}
\usepackage{booktabs}
\usepackage{enumitem}
\usepackage{indentfirst}
\usepackage{tikz}
\usetikzlibrary{positioning,fit,backgrounds}
\tikzset{
  wifiicon/.pic={\fill (0,0) circle (0.05);
    \foreach \r in {0.14,0.25,0.36} \draw[line width=1.1pt] (50:\r) arc (50:130:\r);}}
\colorlet{wifibg}{cyan!25}
\usepackage[hidelinks]{hyperref}


% Dane grupy - uzupełnić
\newcommand{\autorA}{Adam Korwin 284455}
\newcommand{\autorB}{Michał Hałas 277462}

\begin{document}

% ---------------- Strona tytułowa ----------------
\pagenumbering{Alph}
\begin{titlepage}
  \centering
  \includegraphics[height=3.5cm]{pwr_logo.png}\par\vspace{0.6cm}
  {\Large\textsc{Politechnika Wrocławska}\par}
  \vspace{0.2cm}
  {\large Wydział Informatyki i Telekomunikacji\par}
  \vspace{2.5cm}
  {\large Projekt z przedmiotu\par}
  \vspace{0.3cm}
  {\LARGE\bfseries Technologie Sieciowe\par}
  \vspace{1.2cm}
  {\Large Projekt lokalnej sieci komputerowej\\ dla wydawnictwa\par}
  \vspace{0.6cm}
  \vspace{2.5cm}
  \begin{flushleft}
    \textbf{Skład grupy:}
    \begin{enumerate}[leftmargin=2em]
      \item \autorA
      \item \autorB
    \end{enumerate}
  \end{flushleft}
  \vfill
  {\large Wrocław, \today\par}
\end{titlepage}
\pagenumbering{arabic}

\tableofcontents
\newpage

% ---------------- 1. Wstęp ----------------
\section{Wstęp}

\subsection{Charakterystyka przedsiębiorstwa}

Przedmiotem projektu jest sieć komputerowa wydawnictwa z siedzibą we Wrocławiu.
Wydawnictwo prowadzi pełny cykl wydawniczy: od przyjęcia materiałów od autorów
i klientów, przez opracowanie redakcyjne i korektę tekstów, aż po przygotowanie
do druku i druk. Firma zatrudnia 310 pracowników pracujących przy stanowiskach
komputerowych, zorganizowanych w trzy grupy robocze (działy):

\begin{itemize}
  \item \textbf{Drukarnia} (103 stanowiska) -- dział odpowiedzialny za
  przygotowanie publikacji do druku (skład, przygotowanie plików produkcyjnych)
  oraz realizację zleceń drukarskich;
  \item \textbf{Korektorzy} (85 stanowisk) -- dział zajmujący się redakcją
  i korektą tekstów przed ich skierowaniem do druku;
  \item \textbf{Konsultanci} (122 stanowiska) -- dział obsługi autorów
  i klientów: przyjmowanie zamówień, konsultacje oraz kontakt z kontrahentami.
\end{itemize}

Przedsiębiorstwo zajmuje dwa budynki: w Budynku~1 piętra 1--2, a w
Budynku~2 piętra 1--3. Pracownicy wszystkich działów są rozmieszczeni
w obu budynkach. Budynki są oddalone od siebie o 204~m i połączone
światłowodem wielomodowym.
W budynkach zainstalowano okablowanie strukturalne kategorii~6 oraz szafy
teleinformatyczne: główny punkt dystrybucyjny (MDF) na 3.~piętrze Budynku~2
oraz dwa pośrednie punkty dystrybucyjne (IDF1 na 2.~piętrze Budynku~2,
IDF2 na 2.~piętrze Budynku~1). Na 1.~piętrze każdego z budynków pracownicy
korzystają również z urządzeń bezprzewodowych. Rozmieszczenie stanowisk,
urządzeń i punktów dystrybucyjnych przedstawiono na rysunku~\ref{fig:rozmieszczenie}.

\begin{figure}[htbp]
\centering
\begin{tikzpicture}[font=\scriptsize,
  floor/.style={draw, anchor=south west, minimum width=6.9cm, minimum height=2cm,
    text width=6.6cm, align=left, inner sep=4pt},
  wifi/.style={fill=wifibg},
  rack/.style={draw, fill=black!75, text=white, font=\scriptsize\bfseries,
    minimum width=1.1cm, minimum height=0.5cm, inner sep=1pt},
  fiber/.style={very thick, orange!80!black},
  copper/.style={very thick, green!50!black},
  serves/.style={densely dotted, thick, ->, >=stealth}]

  % Budynek 1 (2 piętra), punkt dystrybucyjny w prawym górnym rogu
  \node[floor, wifi] (b1p1) at (0,0)
    {\textbf{Piętro 1} -- 99 stanowisk\\
     Drukarnia: 17, Korektorzy: 43, Konsultanci: 39\\
     Drukarki: 1, AP Wi-Fi: 2, urz. bezprz.: 15};
  \node[floor] (b1p2) at (0,2)
    {\textbf{Piętro 2} -- 64 stanowiska\\
     Drukarnia: 22, Korektorzy: 17, Konsultanci: 25\\
     Drukarki: 3};
  \begin{scope}[on background layer]
    \node[draw, very thick, rounded corners, fill=black!4, inner sep=5pt, fit=(b1p1)(b1p2)] (bud1) {};
  \end{scope}
  \node[font=\small\bfseries, below=2pt of bud1] {Budynek 1};

  % Budynek 2 (3 piętra), punkty dystrybucyjne w prawym górnym rogu
  \node[floor, wifi] (b2p1) at (8.6,0)
    {\textbf{Piętro 1} -- 8 stanowisk\\
     Drukarnia: 5, Korektorzy: 2, Konsultanci: 1\\
     Drukarki: 2, AP Wi-Fi: 1, urz. bezprz.: 12};
  \node[floor] (b2p2) at (8.6,2)
    {\textbf{Piętro 2} -- 77 stanowisk\\
     Drukarnia: 24, Korektorzy: 18, Konsultanci: 35\\
     Drukarki: 2};
  \node[floor] (b2p3) at (8.6,4)
    {\textbf{Piętro 3} -- 62 stanowiska\\
     Drukarnia: 35, Korektorzy: 5, Konsultanci: 22\\
     Drukarki: 2};
  \begin{scope}[on background layer]
    \node[draw, very thick, rounded corners, fill=black!4, inner sep=5pt, fit=(b2p1)(b2p3)] (bud2) {};
  \end{scope}
  \node[font=\small\bfseries, below=2pt of bud2] {Budynek 2};

  % Odległość między budynkami
  \draw[<->, >=stealth, thick] (bud1.east |- 0,1) -- (bud2.west |- 0,1)
    node[midway, above, font=\scriptsize\bfseries] {204 m};

  % Wi-Fi
  \pic[cyan!40!black] at ([shift={(-0.45,0.3)}]b1p1.south east) {wifiicon};
  \pic[cyan!40!black] at ([shift={(-0.45,0.3)}]b2p1.south east) {wifiicon};

  \node[rack, anchor=north east] (idf2) at ([shift={(-0.12,-0.12)}]b1p2.north east) {IDF2};
  \node[rack, anchor=north east] (idf1) at ([shift={(-0.12,-0.12)}]b2p2.north east) {IDF1};
  \node[rack, anchor=north east] (mdf)  at ([shift={(-0.12,-0.12)}]b2p3.north east) {MDF};

  \draw[fiber] (idf2.east) -- ++(0.9,0) |- ([yshift=0.35cm]b2p3.north -| mdf) -- (mdf.north);
  \draw[copper] (idf1.north) -- (mdf.south);
  \draw[serves] ([xshift=0.35cm]idf2.south) -- ([xshift=-0.32cm]b1p1.north east);
  \draw[serves] ([xshift=0.35cm]idf1.south) -- ([xshift=-0.32cm]b2p1.north east);
\end{tikzpicture}

\vspace{6pt}
{\scriptsize
\tikz\draw[very thick, orange!80!black] (0,0) -- (0.6,0); światłowód wielomodowy\quad
\tikz\draw[very thick, green!50!black] (0,0) -- (0.6,0); skrętka kat.~6\quad
\tikz\draw[densely dotted, thick, ->, >=stealth] (0,0) -- (0.6,0); piętro obsługiwane przez IDF\quad
\tikz\draw[fill=wifibg] (0,0) rectangle (0.4,0.22); \tikz[baseline=-1pt]\pic[cyan!40!black,scale=0.8]{wifiicon}; piętro z siecią Wi-Fi\par}
\caption{Rozmieszczenie stanowisk, urządzeń i punktów dystrybucyjnych w budynkach wydawnictwa}
\label{fig:rozmieszczenie}
\end{figure}

Profil działalności wydawnictwa przekłada się bezpośrednio na charakter ruchu
w sieci. Praca na tekstach i plikach przygotowanych do druku wymaga intensywnej
wymiany danych z serwerami lokalnymi oraz z drukarkami sieciowymi, w których
dominuje ruch wysyłany przez użytkowników (zapis dokumentów, zlecenia wydruku).
Firma utrzymuje także własne serwery dostępne z Internetu: serwer WWW
(serwis internetowy wydawnictwa) oraz serwer FTP, służący do wymiany plików
z autorami i klientami. Pracownicy korzystają ponadto z aplikacji
internetowych: przeglądarki, wideokonferencji, telefonii VoIP, klienta FTP,
komunikatora, usług pracy w chmurze oraz poczty elektronicznej. Ciągłość pracy
sieci jest dla firmy istotna, ponieważ jej przerwa zatrzymuje zarówno
produkcję wydawniczą, jak i obsługę klientów.

\subsection{Cel projektu}

Celem projektu jest zaprojektowanie lokalnej sieci komputerowej wydawnictwa,
która zapewni prawidłową i niezawodną obsługę wszystkich pracowników, serwerów,
drukarek oraz urządzeń bezprzewodowych, z uwzględnieniem rzeczywistych
wymagań na przepustowość zgłoszonych przez zamawiającego.

Projektowana sieć ma spełniać następujące wymagania zamawiającego:
\begin{itemize}
  \item wykorzystanie technologii z rodziny Ethernet;
  \item separacja ruchu poszczególnych grup roboczych za pomocą sieci VLAN;
  \item dostęp do sieci bezprzewodowej na wskazanym piętrze każdego budynku;
  \item rezerwa portów na przełącznikach w wysokości 20\,\% portów zajętych,
  wynikająca z przewidywanego wzrostu zatrudnienia;
  \item dwa łącza do Internetu (podstawowe i zapasowe) o przepustowości
  adekwatnej do potrzeb firmy, przy czym łącze podstawowe ma gwarantować
  co najmniej 40\,\% średniego przewidywanego przepływu;
  \item zabezpieczenie sieci przed typowymi zagrożeniami, takimi jak awaria
  zasilania, złośliwe oprogramowanie czy próby włamania z Internetu.
\end{itemize}

\subsection{Zakres projektu}

Okablowanie strukturalne, szafy teleinformatyczne oraz urządzenia końcowe
(komputery, serwery, drukarki) są już zainstalowane i nie wchodzą w zakres
prac. Zakres projektu obejmuje:
\begin{enumerate}
  \item inwentaryzację zasobów sprzętowych, aplikacji i zasobów ludzkich,
  \item analizę potrzeb użytkowników, w tym oszacowanie przepływów danych
  w kluczowych punktach sieci,
  \item sformułowanie założeń projektowych,
  \item projekt logiczny sieci wraz z podziałem na sieci VLAN,
  \item dobór urządzeń sieciowych,
  \item projekt adresacji IP oraz konfiguracji urządzeń,
  \item projekt podłączenia do Internetu wraz z wyborem dostawców,
  \item analizę bezpieczeństwa i niezawodności sieci,
  \item kosztorys obejmujący urządzenia oraz koszt łączy internetowych
  w okresie 2~lat.
\end{enumerate}

% ---------------- Kolejne etapy ----------------
\section{Inwentaryzacja zasobów: sprzętu, aplikacji, zasobów ludzkich}
% Etap 2

\section{Analiza potrzeb użytkowników -- wymagania zamawiającego}
% Etap 2

\section{Założenia projektowe}
% Etap 2

\section{Projekt sieci}

\subsection{Projekt logiczny sieci wraz z opisem koncepcji rozwiązania i uzasadnieniem}
% Etap 3

\subsection{Wybór urządzeń sieciowych}
% Etap 3

\subsection{Projekt adresacji IP}
% Etap 3

\subsection{Projekt konfiguracji urządzeń}
% Etap 4

\subsection{Projekt podłączenia do Internetu}
% Etap 4

\subsection{Analiza bezpieczeństwa i niezawodności sieci}
% Etap 4

\subsection{Kosztorys}
% Etap 4

\section{Karty katalogowe proponowanych urządzeń}
% Etap 4

\end{document}
